\chapter{Conclusions and future work}

\section{Conclusions}
The development of \textit{Micorrizas} has made it possible to create a functional Android serious game based on a real educational activity. The project has achieved its main technical objective by producing a mobile application that combines geolocation, cooperative multiplayer interaction, mission-based gameplay and a final group result. These elements work together to support an outdoor learning experience in the \textit{Jardín de los Sentidos}.

One of the main conclusions of the project is that the game is not simply a gamified version of an existing activity. During development, it became clear that the original visit was being used as the basis for designing a serious game with its own structure, narrative, mechanics and progression. The educational activity provided the context, contents and objectives, but the final result required specific design decisions in order to transform that material into an interactive experience.

The project has also shown the value of connecting mobile gameplay with a real environment. By using the garden as the playable space, the game encourages players to observe, move, communicate and make decisions in relation to the physical surroundings. This supports the objective of promoting awareness and appreciation of biodiversity through interaction with the natural environment.

However, the educational impact of this objective cannot yet be fully measured. The application is prepared to be introduced as a tool in the subject, and the lecturers have expressed the intention of using it in a real teaching session during the next academic year. Once the game is used with students in that context, it will be possible to evaluate more precisely whether the experience increases their awareness, appreciation and understanding of the biodiversity of the \textit{Jardín de los Sentidos}.

Working on mycorrhizae has enabled me to put much of the knowledge I acquired during my degree into practice, as well as to incorporate new knowledge that has proved necessary during the course of the project.

Overall, the result achieved is satisfactory, as the game meets the project’s expectations and demonstrates the feasibility of building a geolocation-based cooperative experience in Unity.


\section{Future work}
The main line of future work is the use of the application in a real session of the subject. This would allow the game to be tested with its intended users and under authentic teaching conditions. The results of that implementation could be used to evaluate the educational impact of the game, especially in relation to student engagement and collaborative learning.

A future evaluation could combine several methods, such as observation during the activity, questionnaires before and after the session, teacher feedback and analysis of the final group scores. This would make it possible to assess not only whether the application works technically, but also whether it contributes to the learning objectives of the original activity.

From a technical perspective, future work could focus on improving GPS accuracy and robustness during outdoor use. Although the prototype already includes geolocation-based mission activation, further testing with larger groups and different mobile devices would help to refine activation areas, GPS indicators and tolerance margins.

The content of the game could also be expanded. New plant species, sounds, images and questions could be added to enrich the missions and make the experience more varied. Additional missions could also be designed for other botanical or educational spaces, allowing the structure of \textit{Micorrizas} to be reused beyond the \textit{Jardín de los Sentidos}.

Finally, the teacher-facing part of the application could be improved. Future versions could include tools to export results, review group performance or configure the content of the missions before the session. These improvements would make the application more useful as a teaching resource and would support its integration into the subject in a more stable way.

\section{Project applications}
\textit{Micorrizas} has been developed as a tool that can be applied within the subject related to the natural environment in Early Childhood Education. Its main application is to complement the existing visit to the \textit{Jardín de los Sentidos} by providing a structured, interactive and cooperative activity that students can complete in the real space. If any problems arise, such as poor lighting because the activity is taking place late in the day or adverse weather conditions.

The game can be used by lecturers to guide students through the garden without requiring constant direct instruction. The application gives players objectives, activates quests according to their position and provides feedback on their answers. This allows the teaching activity to become more autonomous while preserving the importance of direct contact with the natural environment.

Another possible application is its use as an assessable class activity. Since the game records the group’s progression and provides a final score, the lecturers can use the result as part of the evaluation of the session. This does not replace the educational value of observation and discussion, but it provides an additional element that can help structure and assess the activity.

The project also creates a foundation for future educational uses. New plant species, sounds, questions, clues or missions could be added to adapt the experience to future editions of the subject or to different learning objectives. Therefore, the application is not limited to the current application, but can be expanded as a reusable educational resource.

%puedes añadir aquí un apartado





% \chapter{Conclusions and future work}

% \section{Conclusions}

% The development of \textit{Micorrizas} produced a functional Android serious game based on a real educational activity. The prototype combines geolocation, cooperative multiplayer interaction, mission-based gameplay and a final group result in an outdoor learning experience located in the \textit{Jardín de los Sentidos}.

% The original visit provided the context, contents and learning objectives. The final design transformed that material into a serious game with narrative progression, cooperative missions and interaction models adapted to mobile devices.

% The project confirms the feasibility of using the garden as a playable space. Players move through the real environment, observe their surroundings, communicate with the group and make decisions linked to the sensory areas. This design also exposed the main technical challenges of the project: GPS accuracy, outdoor usability, mobile testing and synchronisation between devices.

% The most relevant technical achievement was the integration of several systems into a coherent flow. Multiplayer connection, shared session state, GPS readings, ArcGIS map representation, mission activation and mini-game progression work together through explicit state changes and events.

% The validation process confirmed the Android execution of the prototype and the correct behaviour of the main cooperative flow. It also improved mission clarity, interface readability and the connection between mechanics and educational context. The measurement of learning impact remains pending until the game is used with students in a real class session.

% The project also contributed to my development as a video game developer. It required applying knowledge from the degree and extending it in geolocation, cooperative gameplay, mobile deployment and serious game design.

% Overall, the result provides a solid base for classroom validation, technical refinement and future content expansion.

% \section{Future work}

% The next step is to use the application in a real session of the subject. This would allow the prototype to be tested with its intended users, under outdoor conditions and with teacher participation.

% The evaluation should be organised in stages. A first pilot with a small group of students would help detect usability problems, unclear instructions and technical issues before class deployment.

% The second stage would introduce the game into the real teaching activity in the \textit{Jardín de los Sentidos}. This would provide evidence about group movement, device variability, GPS behaviour and integration with the teacher's workflow.

% A later pedagogical study could include pre-test and post-test questionnaires, observation during the activity, teacher feedback and comparison between groups using the traditional visit and groups using the game as a complementary resource. This would make it possible to evaluate engagement, collaboration and understanding of the garden's biodiversity.

% \section{Technical improvements}

% Future technical work should prioritise GPS robustness and outdoor reliability. Tests with larger groups and different Android devices would help refine activation areas, GPS indicators and tolerance margins.

% The teacher-facing side of the application could include result export, group performance review and mission configuration before the session. These features would make the application easier to integrate into the subject.

% A content editor would also be useful. Lecturers could update plant species, sounds, images, clues or answers without modifying the Unity project directly. This would make the tool easier to maintain across future editions of the activity.

% \section{Content extension}

% The current prototype includes one mission for each sensory area. Future versions could add more plant species, sounds, images and questions to increase variety and replayability.

% New missions should keep a clear educational purpose, a specific interaction model and a direct relationship with the physical environment. This criterion would allow the project to grow while preserving coherence between gameplay and learning objectives.

% \section{Dissemination}

% After classroom validation, the results could be presented as an academic communication. The project connects video game development, educational innovation and environmental education, making it relevant for venues focused on serious games, teaching innovation or educational technology.

% Potential dissemination venues include conferences on video games and interactive media, as well as educational conferences such as EDULEARN or ICERI. A future publication should separate technical feasibility, user experience and pedagogical impact clearly.
